\documentclass[10pt]{article}
\usepackage[T1]{fontenc}
\usepackage[utf8]{inputenc}
\usepackage[french]{babel}
\usepackage{amsmath,amssymb,mathrsfs,mathpazo}
\usepackage{geometry}
\geometry{a4paper,margin=2cm}
\usepackage{graphicx,array,tabularx,tabularray,colortbl,multirow,multicol,booktabs,enumitem,pifont,ulem,fancybox,diagbox,cancel}
\usepackage{pstricks,pstricks-add,pst-plot,pst-tree,pst-node,pst-text,pst-func,pst-math,pst-bspline,pst-3dplot,pst-eucl,pst-solides3d}
\let\pscalculate\undefined
\usepackage{pst-bezier}
\usepackage{tikz,pgfplots,tkz-tab}
\usetikzlibrary{arrows,arrows.meta,patterns,calc,positioning,trees,decorations.markings,shapes.geometric,shapes.misc,shapes.arrows,intersections,decorations.pathmorphing,decorations.pathreplacing,backgrounds,shadows,babel}
\pgfplotsset{compat=1.18}
\usepgfplotslibrary{fillbetween}
\usepackage[np]{numprint}
\usepackage{listings,makecell,pas-tableur,twoopt}
\let\barre\undefined
\usepackage[pstricks]{tablvar}
\definecolor{darkgreen}{RGB}{0,100,0}
\providecommand{\red}{\color{red}}
\providecommand{\blue}{\color{blue}}
\providecommand{\green}{\color{green}}
\providecommand{\black}{\color{black}}
\providecommand{\white}{\color{white}}
\providecommand{\gray}{\color{gray}}
\providecommand{\cyan}{\color{cyan}}
\providecommand{\magenta}{\color{magenta}}
\providecommand{\yellow}{\color{yellow}}
\providecommand{\violet}{\color{violet}}
\providecommand{\orange}{\color{orange}}
\providecommand{\euro}{\texteuro}
\providecommand{\EUR}{\texteuro}
\providecommand{\ieme}{\textsuperscript{e}}
\providecommand{\degres}{\ensuremath{^\circ}}
\providecommand{\degree}{\ensuremath{^\circ}}
\providecommand{\vv}[1]{\overrightarrow{#1}}
\providecommand{\Coord}[1]{#1}
\providecommand{\decosix}{}
\providecommand{\starredbullet}{\ensuremath{\bullet}}
\pagestyle{empty}
\setlength{\parindent}{0pt}
\usepackage{adjustbox}\DeclareUnicodeCharacter{2236}{:}\DeclareUnicodeCharacter{2190}{\ensuremath{\leftarrow}}
\providecommand{\R}{\mathbb{R}}
\providecommand{\N}{\mathbb{N}}
\providecommand{\Z}{\mathbb{Z}}
\providecommand{\Q}{\mathbb{Q}}
\providecommand{\D}{\mathbb{D}}
\providecommand{\C}{\mathbb{C}}
\providecommand{\e}{{\mathrm{e}}}
\providecommand{\vect}[1]{\overrightarrow{#1}}
\providecommand{\vectt}[1]{\overrightarrow{\mathrm{#1}}}
\providecommand{\vvt}[1]{\overrightarrow{\mathrm{#1}}}
\providecommand{\np}[2][]{{#2~#1}}
\providecommand{\vv}[1]{{\overrightarrow{#1}}}
\renewcommand{\d}{{\mathrm{d}}}
\providecommand{\degres}{{^\circ}}
\providecommand{\degre}{{^\circ}}
\providecommand{\blue}{\color{blue}}
\providecommand{\red}{\color{red}}
\providecommand{\enskip}{\quad}
\providecommand{\llbracket}{\lbrack\!\lbrack}
\providecommand{\rrbracket}{\rbrack\!\rbrack}
\providecommand{\euro}{{\text{€}}}
\providecommand{\no}{{\text{n°}}}
\providecommand{\interval}[2]{{[#1;#2]}}
\providecommand{\intervFO}[2]{{[#1;#2[}}
\providecommand{\dotfill}{\ldots}

\providecommand{\ds}{\displaystyle}
\providecommand{\textding}[1]{#1}
\providecommand{\Oij}{\(\mathrm{O};\vec{\imath},\vec{\jmath}\)}
\providecommand{\Oijk}{\(\mathrm{O};\vec{\imath},\vec{\jmath},\vec{k}\)}
\providecommand{\Ouv}{\(\mathrm{O};\vec{u},\vec{v}\)}
\providecommand{\intervFF}[2]{\ensuremath{[#1;#2]}}
\renewcommand{\euro}{\ifmmode\text{\texteuro}\else\texteuro\fi}
\DeclareUnicodeCharacter{00B0}{\ifmmode^{\circ}\else\textdegree\fi}
\DeclareUnicodeCharacter{20AC}{\ifmmode\text{\texteuro}\else\texteuro\fi}
% Énoncés : Éducation nationale. Transcription et corrigés : APMEP, Denis Vergès et contributeurs. Découpage et classement : CoopMaths, Sébastien Lozano, Éric Elter et contributeurs. Adaptation de la présentation : Maths Explorer.
% CC BY-SA : https://coopmaths.fr/www/about/
% https://coopmaths.fr/alea/static/eam/2026/tex/eam_specialite_2026_06_ameriquenord_2_cor.tex
\begin{document}
\noindent\textit{APMEP / CoopMaths. CC BY-SA. Présentation adaptée par Maths Explorer.}\par\medskip
\subsection*{Partie A}
	\begin{center}
		\begin{tikzpicture}
			\begin{axis}[/pgf/number format/.cd, use comma,
				x={6mm}, y={6mm},
				xmin = -1, xmax = 24.5,
				ymin = -0.9, ymax = 10.6,
				xtick distance=1,
				ytick distance=1,
				minor tick num = 1,
				grid=both, axis lines =center,
				grid style={gray!,line width=0.35pt,dotted}]
				\addplot[line width=1.2pt,color=red,smooth,samples=100,domain= 0:24] {8*exp(-0.08*(x-13)*(x-13))};
				\node[below left=1mm] at (axis cs: 0,0) {{\footnotesize 0}};
				\node at(axis cs: 8,2) {$ \mathcal{C}_f $};
				\node[text width=2cm] at(axis cs: 2,9.25) {Puissance électrique (en kW)};
				\node at(axis cs: 21.6,1) {Horaire de la journée};
				\draw[teal,dashed,line width=1pt] (axis cs: 11,-0.5) circle (0.35) (axis cs: 11,-0.15)--(axis cs: 11,5.8)--(axis cs: 0,5.8);
				\draw[teal,dotted,line width=1pt] (axis cs: -0.45,5) circle (0.35) (axis cs: -0.1,5)--(axis cs: 15.5,5)--(axis cs: 15.5,0) (axis cs: 10.5,0) -- (axis cs: 10.5,5);
			\end{axis}
		\end{tikzpicture}
	\end{center}

	
	
	
	
	
	
	
	
	
	
	
	
	
	\begin{enumerate}
		\item À $11$~h, la lecture graphique donne une puissance d'environ $5,75$~kW (à la précision permise par le graphique, voir le tracé en ligne tiretée).

		\item On lit les abscisses des points de la courbe situés de la droite d'équation $y=5$. (voir le tracé en ligne pointillée).
				On obtient approximativement :	\quad $f(x)\geqslant 5$ \quad pour $x\in[10,5 ; 15,5]$.

		\emph{Interprétation }: la puissance des panneaux solaires est supérieure ou égale à $5$~kW entre environ 10~h~30 et 15~h~30, soit pendant environ 5 heures dans la journée.
	\end{enumerate}

	\subsection*{Partie B}

	\begin{enumerate}
		\item Chaque année, le tarif augmente de $6~\%$ : on obtient le tarif d'une année en prenant le tarif de l'année précédente, multiplié par $1,06$.

		La suite $\left(c_{n}\right)$ est donc géométrique, de raison $q=1,06$ (et de premier terme $c_{0}=0,15$).

		\item Comme la suite est géométrique, pour tout entier naturel $n$ : \quad $c_{n}=c_{0}\times q^{\,n}=0,15\times 1,06^{\,n}$.

		\item L'année 2030 correspond à $2020 + 10$, donc $n=10$.

		Le coût pour $1$~kWh en 2030 est donc :\quad $c_{10}=0,15\times 1,06^{10}$.

		\item \begin{enumerate}
			\item \texttt{c} représente le coût (en euros) de $1$~kWh pour l'année considérée ($2020+n$) ;

			\texttt{S} représente l'économie totale cumulée (en euros) réalisée grâce aux panneaux solaires depuis 2020. Pour chaque exécution de la boucle \texttt{while}, on ajoute à la somme économisée précédemment la somme économisée cette année là, en multipliant la quantité d'énergie non achetée (\np[kWh]{2000}), par le prix d'achat. Puis, aux deux lignes suivantes, on met à jour l'indice de l'année, et le prix de l'énergie l'année suivante.

			\item Le programme affiche $16$ : il faut $16$ années pour que l'économie cumulée \texttt{S} dépasse \np{7000}~\euro{}, c'est-à-dire le coût d'installation.

			Selon cette modélisation, l'investissement de Camille sera rentabilisé au bout de $16$ ans, en 2020 + 16 = 2036.
		\end{enumerate}
	\end{enumerate}

\end{document}
